\chapter{主要结论}\label{sec:results}

动态参考点改变了策略所优化的目标：即使两条财富路径停在同一终点，先前经历留下的参考点仍可能使CPT评价不同。每一期需要从有限轨迹中估计当前目标的梯度；时期继续推进时，策略还必须跟上目标的移动。由路径依赖出发，本章先建立冻结目标下的梯度估计结果，再给出连续更新时的投影残差与动态局部遗憾界。

\section{参考点的路径依赖}

\begin{proposition}[相同期末财富与不同的CPT价值]\label{prop:path-dependence}
设$x_{\rm lo}<r_0<x_{\rm f}<x_{\rm hi}$，考虑两条财富路径
\[
 \mathrm{I}: r_0\to x_{\rm hi}\to x_{\rm f},
 \qquad
 \mathrm{II}: r_0\to x_{\rm lo}\to x_{\rm f}.
\]
若
\[
 x_{\rm f}<(1-\etaP)r_0+\etaP x_{\rm hi},
 \qquad
 x_{\rm f}>(1-\etaM)r_0+\etaM x_{\rm lo},
\]
则
\begin{equation}\label{eq:path-reference-gap}
 r_{\rm I}-r_{\rm II}
 =\etaP(1-\etaM)x_{\rm hi}
 -\etaM(1-\etaP)x_{\rm lo}
 +(\etaM-\etaP)x_{\rm f}.
\end{equation}
当右侧不为零时，两条路径虽然具有相同的期末财富$x_{\rm f}$，其正则化CPT价值却不同。
\end{proposition}

这一命题揭示了终值评价与路径评价的区别：投资者在最终持有同样财富时，先前的上涨或下跌仍会通过参考点留下影响。附录~\ref{app:reference-gradient}给出直接递推的证明。

\begin{proposition}[参考点敏感性]\label{prop:reference-sensitivity}
在假设~\ref{ass:trajectory}--\ref{ass:cpt-regularity}下，固定$(P,z,\theta)$。仅在本命题的比较中，假设策略特征和市场转移核与参考点隔离：若使用共同的策略随机数和市场随机数，将期初参考点从$r$改为$r'$不会改变财富与持仓路径。令
\begin{equation}\label{eq:reference-constants}
 \eta_{\min}=\min(\etaP,\etaM),
 \qquad
 L_v=\alpha_+\eps_v^{\alpha_+-1}
 +\lambda\alpha_-\eps_v^{\alpha_--1}.
\end{equation}
沿同一财富路径耦合的两个期初参考点满足
\begin{equation}\label{eq:reference-contraction}
 \abs{r_{t_k+h}(r)-r_{t_k+h}(r')}
 \le(1-\eta_{\min})^h\abs{r-r'},
\end{equation}
并且对固定的$(P,z,\theta)$，
\begin{equation}\label{eq:reference-value-sensitivity}
 \abs{J(P,z,r;\theta)-J(P,z,r';\theta)}
 \le C_1L_v(1-\eta_{\min})^h\abs{r-r'}.
\end{equation}
进一步记
\[
 \Delta_k^{\mathrm{nr}}
 :=\sup_{\theta\in\Theta}
 \abs{J(P_{k+1},z_{k+1},r_k;\theta)-J(P_k,z_k,r_k;\theta)},
\]
则
\begin{equation}\label{eq:variation-decomposition}
 \sup_{\theta\in\Theta}
 \abs{J_{k+1}(\theta,r_{k+1})-J_k(\theta,r_k)}
 \le \Delta_k^{\mathrm{nr}}
 +C_1L_v(1-\eta_{\min})^h\abs{r_{k+1}-r_k}.
\end{equation}
\end{proposition}

式~\eqref{eq:variation-decomposition}把跨期目标变化分为参考点以外的状态、市场变化，以及由期初参考点移动带来的部分。后者并非独立施加的外部扰动，而是账户此前投资路径的结果。

\section{冻结目标的CPT梯度与统计估计}

\begin{theorem}[正则化CPT得分梯度恒等式]\label{thm:gradient-identity}
在假设~\ref{ass:trajectory}--\ref{ass:cpt-regularity}下，固定一个目标，并简记$J=J_k$、$Y=Y_k$、$S_\theta^\pm(y):=S_{P_k,z_k,r_k,\theta}^\pm(y)$、$G_\theta:=G_{k,\theta}$。定义
\begin{align}\label{eq:cpt-multiplier}
 \Psi_\theta(\tau)
 ={}&\int_0^{B_v}(w_+^\eps)'(S_\theta^+(y))
 \{\1[v_+(Y(\tau))>y]-S_\theta^+(y)\}\,dy\notag\\
 &-\int_0^{B_v}(w_-^\eps)'(S_\theta^-(y))
 \{\1[v_-(Y(\tau))>y]-S_\theta^-(y)\}\,dy.
\end{align}
于是
\begin{equation}\label{eq:score-gradient-identity}
 \nabla J(\theta)=\E\{\Psi_\theta(\tau)G_\theta(\tau)\}.
\end{equation}
同一恒等式也适用于冻结模拟模型的目标$\widetilde J$。
\end{theorem}

$\Psi_\theta$根据轨迹结果超过各价值阈值的情况，以及这些阈值在总体中的生存概率，确定该轨迹在CPT评价中的权重；$G_\theta$则说明轨迹出现概率如何随参数改变。两者相乘再取期望，得到目标函数对因子参数的梯度。此后讨论有限样本估计时，固定模型、起点与所查询的参数，记相应冻结目标的梯度为
\[
 g:=\nabla_\theta J(\theta).
\]
在线算法使用的冻结模型为$\widetilde P_k$，因此它在给定第$k$期期初信息后的估计目标是$\widetilde g_k(\theta_k)$。

若直接把经验生存概率代入非线性的权重导数，有限样本梯度通常带有偏差。附录~\ref{app:estimator}证明，留一法基项的期望误差为$O(n^{-1})$；全样本与两个半样本共享轨迹后，一阶影响项相互抵消，层间差分的二阶矩降为$O(n^{-2}4^{-\ell})$。算法~\ref{alg:cloo-rd}的随机层修正据此把有限精度基项连接到无限精度目标。

\begin{theorem}[严格无偏、有限方差与有限期望成本]\label{thm:estimator}
令$Z_n^{\db}$由式~\eqref{eq:debiased-output}定义，其中$n\ge2$、$p_{\rm act}\in(0,1]$、$1<b<2$。令$C_U$和$C_{\rm rem}$为引理~\ref{lem:level-cancellation}中的有限常数，记
\begin{equation}\label{eq:level-series-constants}
 \Xi_b
 :=\sum_{\ell\ge1}\frac{4^{-\ell}}{\nu_\ell}
 =\frac{1}{4(1-2^{-b})(1-2^{b-2})},
 \qquad
 \Upsilon_b
 :=\sum_{\ell\ge1}\nu_\ell2^\ell
 =\frac{2(1-2^{-b})}{1-2^{1-b}}.
\end{equation}
则
\begin{equation}\label{eq:estimator-unbiased}
 \E Z_n^{\db}=g,
\end{equation}
\begin{equation}\label{eq:single-output-variance}
 \tr\Cov(Z_n^{\db})
 \le C_{\rm var}(n,p_{\rm act},b)
 :=\frac{C_U}{n}+\frac{9C_{\rm rem}\Xi_b}{p_{\rm act} n^2}<\infty,
\end{equation}
且每次输出使用的完整轨迹数的期望为
\begin{equation}\label{eq:expected-estimator-cost}
 C_{\rm cost}(n,p_{\rm act},b)
 =n(1+p_{\rm act}\Upsilon_b)<\infty.
\end{equation}
对$M$个独立输出，
\begin{equation}\label{eq:estimator-mse}
 \widehat g_{M,n}^{\db}:=\frac1M\sum_{m=1}^{M}Z_{n,m}^{\db},
 \qquad
 \E\norm{\widehat g_{M,n}^{\db}-g}_2^2
 \le\frac{C_{\rm var}(n,p_{\rm act},b)}{M}.
\end{equation}
\end{theorem}

此处的无偏对象是固定模型和期初状态后的CPT梯度。层级分布不截断，保证随机修正的期望覆盖全部层级；指数条件$1<b<2$同时使期望轨迹成本和输出方差有限。

\begin{theorem}[固定目标的渐近正态性]\label{thm:clt}
在定理~\ref{thm:estimator}的条件下，固定$(n,p_{\rm act},b)$，记$\Sigma_n=\Cov(Z_n^{\db})$。则
\begin{equation}\label{eq:clt}
 \sqrt M\,(\widehat g_{M,n}^{\db}-g)
 \rightsquigarrow \mathcal N(0,\Sigma_n).
\end{equation}
若$\Sigma_n$非奇异，$\widehat\Sigma_{M,n}$为$M$次输出的经验协方差，则
\begin{equation}\label{eq:studentized-clt}
 \sqrt M\,\widehat\Sigma_{M,n}^{-1/2}
 (\widehat g_{M,n}^{\db}-g)
 \rightsquigarrow \mathcal N(0,\mathrm{Id}_d).
\end{equation}
\end{theorem}

独立重复输出因而能够用于估计固定一期梯度的不确定性。跨期目标即使继续移动，也可以在每一期分别使用这一统计结论。

\section{固定维度下的极小极大信息率}

\begin{definition}[完整轨迹查询实验]\label{def:minimax-oracle}
固定维度$d$、$\theta_0=0\in\operatorname{int}(\Theta)$、$h$和$N$的共同有限上界，以及假设~\ref{ass:trajectory}--\ref{ass:cpt-regularity}中的常数。设模型类$\mathfrak C$满足这些界，包含附录~\ref{app:minimax}构造的伯努利收益子模型，并有$(w_+^\eps)'(0)>0$。对$P\in\mathfrak C$，第$t$次自适应查询$\theta_t\in\Theta$返回一条新的完整增广轨迹
\[
 \tau_t\mid\mathcal G_{t-1}\sim P_{\theta_t},
 \qquad
 \mathcal G_t:=\sigma\{(\theta_j,\tau_j):j\le t\}.
\]
可行程序选择停止时刻$T_{\rm call}$和关于$\mathcal G_{T_{\rm call}}$可测的估计量$\widehat g$，满足
\begin{equation}\label{eq:minimax-budget}
 \sup_{P\in\mathfrak C}\E_P T_{\rm call}\le\mathsf B,
\end{equation}
估计目标为$g_P:=\nabla J_P(\theta_0)$。
\end{definition}

相应的极小极大风险为
\begin{equation}\label{eq:minimax-risk}
 \mathfrak R_{\mathsf B}
 :=\inf_{\widehat g:\,\sup_{P\in\mathfrak C}\E_P T_{\rm call}\le\mathsf B}
 \sup_{P\in\mathfrak C}
 \E_P\norm{\widehat g-g_P}_2^2.
\end{equation}

\begin{theorem}[固定维度下的最优轨迹预算率]\label{thm:minimax}
对定义~\ref{def:minimax-oracle}的查询实验，存在与$\mathsf B$无关的常数$c_{\rm lb},C_{\rm ub},\mathsf B_0\in(0,\infty)$，使任意$\mathsf B\ge\mathsf B_0$均满足
\begin{equation}\label{eq:minimax-rate}
 \frac{c_{\rm lb}}{\mathsf B}
 \le \mathfrak R_{\mathsf B}
 \le \frac{C_{\rm ub}}{\mathsf B}.
\end{equation}
在相同期望完整轨迹预算下，独立平均未截断的$Z_n^{\db}$输出可以达到该上界的常数倍。因此，当$d$和正则化参数固定时，梯度估计的极小极大速率为$\Theta(\mathsf B^{-1})$。
\end{theorem}

\section{移动CPT目标下的一阶跟踪}

\begin{theorem}[期望平均平方投影残差上界]\label{thm:online-residual}
设假设~\ref{ass:trajectory}--\ref{ass:cpt-regularity}成立，且$\gamma L\le\vartheta$。固定估计器设计$(n,p_{\rm act},b)$，令$C_g:=\max\{B_g,\vartheta\}$，$\widehat g_k^{\db}$由式~\eqref{eq:gradient-average}给出，目标为$\widetilde g_k(\theta_k)$。则
\begin{equation}\label{eq:true-gradient-mse}
 \E\norm{\widehat g_k^{\db}-g_k(\theta_k)}_2^2
 \le\frac{C_{\rm var}(n,p_{\rm act},b)}{M_k}+\E\beta_k^2,
\end{equation}
并有
\begin{equation}\label{eq:online-residual-bound}
 \boxed{
 \cE_K
 \le
 \frac{8\{2B_v+\E\cV_K\}}{\vartheta\gamma K}
 +\frac{10C_g}{\vartheta^3K}
 \sum_{k=1}^{K}
 \left(\frac{C_{\rm var}(n,p_{\rm act},b)}{M_k}+\E\beta_k^2\right).
 }
\end{equation}
\end{theorem}

上界中的$\cV_K$对应真实目标的移动，$M_k^{-1}$对应蒙特卡洛估计误差，$\beta_k^2$对应冻结模拟模型与真实模型之间的梯度差异。控制的是每一期当下的一阶条件，而不是要求所有时期共享同一个最优参数。

\begin{corollary}[动态局部遗憾]\label{cor:dynamic-local-regret}
任取固定窗口长度$W\ge1$、折扣系数$\rho\in(0,1]$及总期数$K\ge1$，逐路径有
\begin{equation}\label{eq:dynamic-local-regret-bound}
 \Reg_{\rho,W}(K)
 \le \sum_{k=1}^{K}\zeta_k,
\end{equation}
从而
\begin{equation}\label{eq:average-dynamic-local-regret-bound}
 \frac1K\E\Reg_{\rho,W}(K)
 \le \cE_K.
\end{equation}
因此，投影残差界也控制了对近期一阶残差作时间平滑后得到的动态局部遗憾。
\end{corollary}

投影残差直接衡量受约束的一阶驻点条件。若还满足统一的局部误差界：存在$\kappa<\infty$使
\begin{equation}\label{eq:error-bound}
 \dist(\theta,\Sta_k)
 \le \kappa\norm{Q_k(\theta)}_2,
 \qquad \theta\in\Theta,
\end{equation}
则残差大小进一步给出参数到当期驻点集合的距离。

\begin{corollary}[移动驻点集合的跟踪]\label{cor:stationary-tracking}
若式~\eqref{eq:error-bound}成立，则
\begin{equation}\label{eq:stationary-tracking}
 \frac1K\sum_{k=1}^{K}
 \E\dist^2(\theta_k,\Sta_k)
 \le\kappa^2\cE_K.
\end{equation}
\end{corollary}

\begin{corollary}[平均平方投影残差与动态局部遗憾趋于零]\label{cor:vanishing-tracking}
固定$W<\infty$和$\rho\in(0,1]$。若
\begin{equation}\label{eq:vanishing-conditions}
 \E\cV_K=o(K),
 \qquad
 \frac1K\sum_{k=1}^{K}
 \left(\frac1{M_k}+\E\beta_k^2\right)\to0,
\end{equation}
则
\begin{equation}\label{eq:vanishing-conclusion}
 \cE_K\to0,
 \qquad
 K^{-1}\E\Reg_{\rho,W}(K)\to0.
\end{equation}
若式~\eqref{eq:error-bound}也成立，则式~\eqref{eq:stationary-tracking}的左侧同样趋于零。
\end{corollary}
